\documentclass{article}
\usepackage{graphicx} % Required for inserting images


\begin{document}

II. Si un objeto está a una distancia \( 2f \) de un espejo cóncavo con foco \( f \), ¿cuál es la posición de su imagen?  
  
La posición de la imagen en un espejo cóncavo se determina usando la ecuación del espejo esférico:

\[
\frac{1}{f} = \frac{1}{s} + \frac{1}{s_i}
\]

Dado que el objeto está a \( s = 2f \), sustituimos en la ecuaciómn:  

\[
\frac{1}{f} = \frac{1}{2f} + \frac{1}{s_i}
\]

Despejando \( s_i \):  

\[
\frac{1}{s_i} = \frac{1}{f} - \frac{1}{2f} = \frac{2}{2f} - \frac{1}{2f} = \frac{1}{2f}
\]

\[
s_i = 2f
\]

La imagen se forma en \( s_i = 2f \), es decir, a la misma distancia que el objeto. Además, la imagen es real, invertida y de igual tamaño que el objeto. (con la fórmula de Newton que no vimos en el pdf de lab esto queda directo pues $x_0x_i=f^2 =x_0f=(s_i-f)f \Rightarrow s_i=f+f=2f$)

III. ¿Qué implicación tiene la aproximación paraxial en el montaje de un experimento?  
 
La aproximación paraxial se basa en considerar únicamente los rayos que forman un ángulo pequeño con el eje óptico y que están cerca de este. Esto implica que se pueden usar las aproximaciones físicas misteriosas Tayloríficas (McLaurin) \( \sin \theta \approx \theta \) y \( \tan \theta \approx \theta \) para simplificar los cálculos de trayectorias de los rayos.  

En un experimento de óptica, la aproximación paraxial tiene las siguientes implicaciones:  
1. Reducción de aberraciones: Al trabajar con rayos cercanos al eje óptico, se desprecian las aberraciones esféricas y otras distorsiones ópticas. Además es especialmente conveniente para obtener triágulos "rectángulos" semejantes y formular teoría geométrica  
2. Uso de fórmulas simplificadas: Permite aplicar las ecuaciones de los espejos y lentes delgadas sin correcciones complejas.  
3. Precisión en la formación de imágenes: Dentro del régimen paraxial, los resultados obtenidos son suficientemente precisos para muchos experimentos ópticos sin necesidad de recurrir a óptica geométrica avanzada o trazado exacto de rayos.  

La aproximación paraxial permite simplificar los cálculos ópticos, minimizar aberraciones y mejorar la precisión en la formación de imágenes en experimentos ópticos.

\end{document}
